The interactor first prints an integer $T$, the number of test cases.

For each test case, Minh initially stands at the root of the tree.

To move through the cave, you may print one of the following queries:

\begin{itemize}
\item $\texttt{? L}$ --- attempt to move from the current chamber to its left child;
\item $\texttt{? R}$ --- attempt to move from the current chamber to its right child;
\item $\texttt{? B}$ --- attempt to move from the current chamber to its parent.
\end{itemize}

After each query, the interactor prints an integer:

\begin{itemize}
\item $1$ if the requested move is possible. In this case, Minh moves to the corresponding chamber;
\item $0$ if the requested move is impossible. In this case, Minh remains in the current chamber.
\end{itemize}

In particular, the query $\texttt{? B}$ returns $0$ when Minh is currently at the root.

You may perform at most $8000$ queries for each test case.

When you have determined the total number of chambers, print

$$\texttt{! N}$$

where $N$ is the total number of chambers in the current test case.

Printing the final answer does not count as a query.

If the answer is correct, the interaction proceeds to the next test case. Otherwise, your submission receives a Wrong Answer verdict.

After printing every query, do not forget to flush the output. You must also flush the output after printing the final answer.